%!TEX root = ../main.tex
%
% イントロダクション
%


\section{はじめに}

\subsection{自然科学とは}

\begin{itemize}

\item 科学とは、論理的に物事を理解しようとする試みです。したがって、すべての学問は科学であるとも言えます。

\item 自然科学とは、自然現象を対象とした学問です。したがって、数学は自然科学ではありません。なぜなら数学は自然を対象としていないからです。

\item 自然科学は自然を理解し、その自然現象が起こる原理(原因と結果)を統一的に解明することを目的とします。いかに優れた理論であっても自然を正しく反映していなければそれは正しい理論ではありません。

\item 理論は実験観測によって確認されなければなりません。(演繹的アプローチ)

\item 理論を作るためには実験観測の結果を分析しなければなりません。 (帰納的アプローチ)

\end{itemize}

\subsubsection*{【調べてみよう】}

アインシュタインは相対性理論ではノーベル賞をもらえませんでした。何故でしょうか?

\subsection{科学における数学}

\begin{itemize}

\item 数学は土地の測量(幾何学)、物体の運動の解明(微分積分学)のように、応用を目的として発展したものが多いのですが、応用・実用を目的としているわけではありません。

\item 多くの学問では、数学は論理的に考え理解する手段として使われます。数学による表現は言語による表現よりあいまいさがありません。数量的に厳密です。数学は人文科学、社会科学を含むすべての科学に利用されています。

\item 数学(数字)は正しく理解して使う必要があります。

\item 物理でよく使う数学は微分、積分、統計学、線形代数です。この方法論では高度な数学は特に必要ありません。

\end{itemize}


\subsection{実験・測定}

\subsubsection*{【考えてみよう】}

\begin{enumerate}

\item 理科年表によると水素の密度は0.0899 kg/m${}^3$,アルゴンの密度は1.784 kg/m${}^3$です。どちらがよい精度で(より正確に)分かっているでしょうか?

\item アルゴンの密度を1.7840 kg/m${}^3$と書くのと上のように1.784 kg/m${}^3$と書くのは何が違うでしょうか?

\item  6人の学生がある物体の重さを量ったら次のようになりました。(単位kg)

\begin{tabular}{|c|l|l|}
\hline
人&測定値&誤差\\
\hline
A&3.5&0.4\\
B&2.94&0.15\\
C&3.27&0.38\\
D&4.00&0.75\\
E&2.85&0.36\\
F&2.99&0.13\\
\hline
\end{tabular}
\qquad 平均値=3.258 kg

この物体の重さはいくらと考えられますか? 平均値をもって、物体の 
重さとみなしてよいでしょうか? 特異な測定値4.00$\pm$0.75 kgはどのよう 
に扱ったらよいでしょうか?

\item ある長方形の土地の面積を求めます。各辺の長さを測定したところ、縦4.0$\pm$0.08 m, 横6.0$\pm$0.12 mでした。土地の面積は24 m${}^2$ですが、測定による面積の誤差はどのくらいでしょうか?

\end{enumerate}

\begin{itemize}

\item 測定値と誤差

\begin{itemize}

\item 物理量は次元を持ち、単位をつけて表します。

\item 測定値には必ず誤差があります。

\item 誤差の原因。誤差は丸め誤差(四捨五入による誤差)だけではありません。統計誤差と系統誤差もあります。

\item 測定値の表し方、有効数字\\
測定値には原則として誤差をつけて示します。誤差をつけないで示されている場合は、最後の桁の下の桁に誤差があると理解します。

\item 測定値の正規分布

\item 統計誤差を減らすためには同じ測定を多数繰り返します。

\end{itemize}


\item 真の値と最確値

\begin{itemize}

\item 真の値は誰も知りません。

\item 複数の測定値からどのようにしてもっとも真の値に近いと思われる値(最確値)を求めることができるでしょうか?

\begin{itemize}
\item 平均値\\
測定値が平均値の周りに集中している場合に使える。測定値がたくさんあるときにのみ使える。
\item 最小二乗法(実験の中で勉強)

\end{itemize}

\end{itemize}

\item 測定値同士の掛け算、割り算などでは有効数字を超えた桁数は無意 
味です。\\
例　10.0 kgのつぼに0.012 kgの砂糖をいれたら、全体の重さはいくらと考えられますか? \\
例　レーザー光を使って２地点AとBの間の距離を測りました。レーザー光が地点Bで反射して戻ってくるまでにかかる時間を測定したところ0.00014 秒でした。AB間の距離は
\[
299792458 \text{ [m/s]}\times 0.00014\text{ [s]} \times \frac{1}{2} 
= 20985.47206 \text{ [m]}
\]
と考えていいでしょうか? (ヒント：測定値から、光が往復するのに 
かかった時間は何秒から何秒の間にあると考えられますか?)

\item 誤差の二乗伝播\\
測定値が$X\pm\delta x$, $Y\pm\delta y$なら\\
長さの和：$X+Y\pm\sqrt{(\delta x)^2+(\delta y)^2}$\\
面積：$XY\left(1\pm \sqrt{(\frac{\delta x}{X})^2+(\frac{\delta y}{Y})^2}\right)$

\subsubsection*{問題}
長さの和を$X+Y\pm (\delta x+\delta y)$とするとなぜよくないのでしょうか? (ヒント：$\delta x +\delta y$と$\sqrt{(\delta x)^2+(\delta y)^2}$はどちらのほうが大きいですか?)

\end{itemize}


\subsection{便利な数学}

\subsubsection*{【指数表示】}

\begin{tabular}{rcl}
電子の質量 &=& 0.000 000 000 000 000 000 000 000 000 000 910 938 370 15 kg\\
&=& 9.109 383 7015$\times 10^{-31}$ kg\\
アボガドロ定数&=& 602 214 076 000 000 000 000 000 mol${}^{-1}$\\
&=&6.022 140 76$\times 10^{23}$ mol${}^{-1}$\\ 
1モルの電子の重さ&=&$9.1093837015\times 10^{-31}\times 6.02214076 \times 10^{23}$\\
&=&$(9.1093837015\times6.02214076) (10^{-31}\times10^{23})$\\
&=&$5.48579909\times10^{-7}$ kg 
\end{tabular}

指数部と仮数部を別々に計算できます。

有効数字を明確に表示できます。(1240 ｍの有効数字はどこまであるか不明確です。$1.24\times10^3$ mとあらわせば、最後の0は有効数字ではないことを示すことができます。)

\subsubsection*{【SI接頭語】}

一兆くらいまでの数は3桁ごとに接頭語をつけて呼ぶことに決めてあります。たとえば2億5千6百万(256,000,000)バイトと言わないで、256メガバイトと呼びます。この接頭語を記号で表し、メガはMで表します。我々の実験ではmA、$\mu$A、kHzなどが登場します。

\subsubsection*{【調べてみよう】}

\begin{itemize}

\item 身近な容器(ペットボトル、ジュース缶、牛乳パック、化粧水、等)に表記してある重さ、体積はどのような単位でかかれていますか?

\item SI接頭語(特にギガからナノまで)はどのようなところでどのように使われていますか? (日本ではデカは使われていません。)

\end{itemize}

\begin{tabular}{|l|c|l|l|c|l|}
\hline
大きさ&記号&呼び方&大きさ&記号&呼び方\\
\hline
$10^{24}$&Y&ヨッタ&$10^{-1}$&d&デシ\\
$10^{21}$&Z&ゼッタ&$10^{-2}$&c&センチ\\
$10^{18}$&E&エクサ&$10^{-3}$&m&ミリ\\
$10^{15}$&P&ペタ&$10^{-6}$&$\mu$&マイクロ\\
$10^{12}$&T&テラ&$10^{-9}$&n&ナノ\\
$10^{9}$&G&ギガ&$10^{-12}$&p&ピコ\\
$10^{6}$&M&メガ&$10^{-15}$&f&フェムト\\
$10^{3}$&k&キロ&$10^{-18}$&a&アト\\
$10^{2}$&h&ヘクト&$10^{-21}$&z&ゼプト\\
$10$&da&デカ&$10^{-24}$&y&ヨクト\\
\hline
\end{tabular}\\
\quad SI 接頭語

\vspace*{1cm}

\begin{tabular}{|c|c|l|c|c|l|}
\hline
大文字&小文字&呼び方&大文字&小文字&呼び方\\
\hline
A&$\alpha$&アルファ&N&$\nu$&ニュー\\
B&$\beta$&ベータ&$\Xi$&$\xi$&グザイ、クシー\\
$\Gamma$&$\gamma$&ガンマ&O&o&オクミロン\\
$\Delta$&$\delta$&デルタ&$\Pi$&$\pi$&パイ\\
E&$\epsilon$&イプシロン&$P$&$\rho$&ロー\\
Z&$\zeta$&ツェータ&$\Sigma$&$\sigma$&シグマ\\
H&$\eta$&エータ&T&$\tau$&タウ\\
$\Theta$&$\theta$&テータ、シータ&$\Upsilon$&$\upsilon$&ウプシロン\\
I&$\iota$&イオタ&$\Phi$&$\phi$&ファイ\\
K&$\kappa$&カッパ&X&$\chi$&カイ\\
$\Lambda$&$\lambda$&ラムダ&$\Psi$&$\psi$&プサイ、プシー\\
M&$\mu$&ミュー&$\Omega$&$\omega$&オメガ\\
\hline
\end{tabular}\\
\quad ギリシャ文字 

\newpage


\subsection{実験レポートの書き方}

実験が終わったら実験の結果を整理し、科学的な考察を行うためにレポートを作成します。レポートの形式は実験の内容によって異なりますが、次にあげる要件を最低限満たすようにします。\\
\begin{minipage}[t]{9.5cm}
\begin{enumerate}
\item
レポートは\underline{A4のレポート用紙}で作成しましょう。(ルーズリーフは使用しないこと。) ただし、レポ ートの作成にパソコンを使用する場合は、レポート用紙ではなくA4 の紙に印刷したものでかまいません。必ず\underline{左上}をホッチキスで留めて提出するようにしましょう。

\item レポートの表紙には、実験の題目、実験日、自分の学籍番号と名前を記入し、その下に共同実験者(同じグループで実験を行った者)の学籍番号と名前を書きましょう。
\end{enumerate}
\end{minipage}
\hspace*{1cm}
\begin{minipage}[t]{5cm}
\begin{center}
\vspace*{-0.3cm}
\includegraphics[scale=0.4]{00_Intro/report.eps}\\
レポートの表紙の例
\end{center}
\end{minipage}

\begin{enumerate}
\setcounter{enumi}{2}
\item レポートは、次のような順序で書くようにしましょう。

\begin{enumerate}

\item 実験の目的:
\begin{itemize}
\item 何を目的とした実験なのかを、簡潔に書くこと。
\end{itemize}

\item 実験の理論の概要

\item 実験装置、実験の手順の概要

\item 実験結果:
\begin{itemize}
\item まず、何を測定して、数値がどうであったかを書き、次に、計算の結果求められた物理量を(途中の式も含めて)書くこと。
\item 数値を書くときには、\underline{単位}を必ず記入すること。
\end{itemize}

\item 検討・考察:

\begin{itemize}
\item 実験結果の誤差について検討するなど、理論的考察を行なう。この実験から
わかったことをまとめ、実験をしてみて新たに感じた疑問があれば、そのよ
うな疑問についても調べ、考察する。
\item 調べて分かったことを記入した際には、必ず参考文献(本の題名、著者、出
版社)を最後に書いておくこと。
\end{itemize}


\end{enumerate}

\end{enumerate}

